\documentclass[12pt,a4paper]{report}
\usepackage[utf8]{inputenc}
\usepackage[T1]{fontenc}
\usepackage[french]{babel}
\usepackage{amsmath,amssymb,amsfonts}
\usepackage{geometry}
\usepackage{physics}

\geometry{hmargin=2.5cm,vmargin=3cm}

\title{\textbf{Calcul des champs électromagnétiques dans les guides d'onde circulaires et leurs discontinuités}}
\author{Théorie des Guides d'Onde \& Discontinuités}
\date{\today}

\begin{document}

\maketitle

\tableofcontents
\newpage

% ==========================================================
\chapter{Calcul des champs électromagnétiques dans un guide d'onde circulaire}
% ==========================================================

\section{Introduction et équations de Helmholtz}

Partant des équations de Maxwell:

\begin{align}
\vec{\nabla} \cdot \vec{E} &= \frac{\rho}{\varepsilon_0} && \text{(Loi de Gauss)} \\
\vec{\nabla} \cdot \vec{B} &= 0 && \text{(Loi de Gauss pour le magnétisme)} \\
\vec{\nabla} \times \vec{E} &= -\frac{\partial \vec{B}}{\partial t} && \text{(Loi d'Faraday)} \\
\vec{\nabla} \times \vec{B} &= \mu_0 \vec{j} + \mu_0 \varepsilon_0 \frac{\partial \vec{E}}{\partial t} && \text{(Loi d'Ampère-Maxwell)}
\end{align}

Considérons un guide d'onde cylindrique parfait de rayon $a$, rempli d'un milieu linéaire, homogène et isotrope caractérisé par sa permittivité $\varepsilon$ et sa perméabilité $\mu$. On adopte le système de coordonnées cylindriques $(r, \theta, z)$, avec l'axe $z$ orienté selon la direction de propagation.
en considérant un courant nul: $\vec{j}=0$ et une densité de charges nulle: $\rho=0$ dans le guide, et on se place en régime harmonique temporel (composante temporelle en $e^{i\omega t}$), en substituant dans les équations de Maxwell et en introduisant

\begin{equation}
\vec{B} = \mu \vec{H}
\end{equation}

Les équation deviennent:

\begin{align}
\vec{\nabla} \cdot \vec{E} &= 0 \\
\vec{\nabla} \cdot \vec{H} &= 0 \\
\vec{\nabla} \times \vec{E} &= -i \omega \mu \vec{H} \\
\vec{\nabla} \times \vec{H} &= i \omega \varepsilon \vec{E}
\end{align}

En combinant les termes $\vec{\nabla} \times \vec{E}$ et $\vec{\nabla} \times \vec{H}$  des deux dernieres équations, on obtient les équations de Helmholtz pour les composantes longitudinales $E_z$ et $H_z$ :
\begin{equation}
\left( \nabla^2 + k^2 \right) \begin{pmatrix} E_z \\ H_z \end{pmatrix} = 0
\end{equation}
où $k = \omega \sqrt{\varepsilon\mu}$ est le nombre d'onde dans le milieu.

En supposant une propagation selon $z$ sous la forme $e^{-i \delta_z \beta z}$, où $\delta_z$ représente le sens de propagation de l'onde suivant z (+1 ou -1), l'opérateur laplacien en coordonnées cylindriques appliqué aux composantes longitudinales donne :
\begin{equation}
\frac{\partial^2 \psi}{\partial r^2} + \frac{1}{r}\frac{\partial \psi}{\partial r} + \frac{1}{r^2}\frac{\partial^2 \psi}{\partial \theta^2} + k_c^2 \psi = 0
\label{eq:helmholtz_cyl}
\end{equation}
avec $\psi \in \{E_z, H_z\}$ et où $k_c^2 = k^2 - \beta^2$ définit la constante de propagation transversale (fréquence de coupure).

\section{Séparation des variables et fonctions de Bessel}
Par la méthode de séparation des variables, on pose $\psi(r, \theta) = R(r)\Phi(\theta)$. L'équation \eqref{eq:helmholtz_cyl} se décompose en deux équations différentielles ordinaires :
\begin{align}
\frac{d^2\Phi}{d\theta^2} + n^2 \Phi &= 0 \quad (n \in \mathbb{N}) \\
r^2 \frac{d^2R}{dr^2} + r \frac{dR}{dr} + \left( (k_c r)^2 - n^2 \right) R &= 0
\end{align}

La solution angulaire impose une périodicité $2\pi$, d'où $\Phi(\theta) = A \cos(n\theta) + B \sin(n\theta)$ avec $n$ entier.
L'équation radiale est l'équation différentielle de Bessel. Sa solution générale s'écrit :
\begin{equation}
R(r) = C J_n(k_c r) + D Y_n(k_c r)
\end{equation}
où $J_n$ est la fonction de Bessel de première espèce d'ordre $n$, et $Y_n$ la fonction de Neumann. Le champ devant rester fini en $r = 0$, la constante $D$ doit être nulle car $\lim_{r\to 0} Y_n(k_c r) = -\infty$.

Ainsi, les solutions longitudinales sont de la forme :
\begin{equation}
\psi(r,\theta,z) = C_n J_n(k_c r) \cos(n\theta) e^{-i\delta_z\beta z}
\end{equation}

\section{Calcul des composantes transversales}
À partir des équations de Maxwell, les composantes transversales se déduisent directement des composantes longitudinales $E_z$ et $H_z$ (en ré-injectant les solutions dans les équations de Maxwell $\vec{\nabla} \times \vec{E}$ et $\vec{\nabla} \times \vec{H}$:
\begin{align}
E_r &= \frac{-i}{k_c^2} \left( \delta_z \beta \frac{\partial E_z}{\partial r} + \frac{\omega\mu}{r} \frac{\partial H_z}{\partial \theta} \right) \\
E_\theta &= \frac{-i}{k_c^2} \left( \delta_z \frac{\beta}{r} \frac{\partial E_z}{\partial \theta} - \omega\mu \frac{\partial H_z}{\partial r} \right) \\
H_r &= \frac{-i}{k_c^2} \left( \delta_z \beta \frac{\partial H_z}{\partial r} - \frac{\omega\varepsilon}{r} \frac{\partial E_z}{\partial \theta} \right) \\
H_\theta &= \frac{-i}{k_c^2} \left( \delta_z \frac{\beta}{r} \frac{\partial H_z}{\partial \theta} + \omega\varepsilon \frac{\partial E_z}{\partial r} \right)
\end{align}

\section{Conditions aux limites et modes $TM$ et $TE$}

\subsection{Modes Transversaux Magnétiques ($TM_{np}$ : $H_z = 0$)}
Pour les modes $TM$, le champ électrique longitudinal $E_z$ doit s'annuler sur la paroi métallique conductrice parfaite ($r = a$) :
\begin{equation}
E_z(a, \theta, z) = 0 \implies J_n(k_c a) = 0
\end{equation}
Soit $x_{np}$ la $p$-ième racine de la fonction de Bessel $J_n(x)$. La constante de coupure vaut :
\begin{equation}
k_{c,np}^{TM} = \frac{x_{np}}{a}
\end{equation}

\subsection{Modes Transversaux Électriques ($TE_{np}$ : $E_z = 0$)}
Pour les modes $TE$, la condition $E_\theta(a, \theta, z) = 0$ impose :
\begin{equation}
\left. \frac{\partial H_z}{\partial r} \right|_{r=a} = 0 \implies J'_n(k_c a) = 0
\end{equation}
Soit $x'_{np}$ la $p$-ième racine de la dérivée de la fonction de Bessel $J'_n(x)$. La constante de coupure vaut :
\begin{equation}
k_{c,np}^{TE} = \frac{x'_{np}}{a}
\end{equation}

\subsection{Récapitulatif: expression finale des composantes des modes circulaires}
La constante de propagation guidée s'obtient pour chaque mode par :
\begin{equation}
\beta_{np} = \sqrt{k^2 - (k_{c,np})^2}
\end{equation}

\paragraph{modes TE}
\begin{align}
E_r &= \frac{i \mu \omega}{k_c ^ 2} \frac{n}{r} J_n(r k_c)\sin(n \theta) \\
E_\theta &= \frac{i \mu \omega}{k_c ^ 2} k_c J'_n(r k_c)\cos(n \theta) \\
H_r &= - \delta_z \frac{i \beta}{k_c ^ 2} k_c J'_n(r k_c)\cos(n \theta) \\
H_\theta &= \delta_z \frac{i \beta}{k_c ^ 2} \frac{n}{r} J_n(r k_c)\sin(n \theta)
\end{align}

les champs dans le plan transverse s'écrivent alors:

\begin{align}
\vec{E_t} = \frac{i \mu \omega}{k_c ^ 2} \begin{pmatrix} \frac{n}{r} J_n(r k_c)\sin(n \theta) \\ k_c J'_n(r k_c)\cos(n \theta) \end{pmatrix} e^{-i\delta_z\beta z} \\
\vec{H_t} = \delta_z \frac{i \beta}{k_c ^ 2} \begin{pmatrix} -k_c J'_n(r k_c)\cos(n \theta) \\ \frac{n}{r} J_n(r k_c)\sin(n \theta) \end{pmatrix} e^{-i\delta_z\beta z}
\end{align}

\paragraph{modes TM}
\begin{align}
E_r &= - \delta_z \frac{i \beta}{k_c ^ 2} k_c J'_n(r k_c)\cos(n \theta) \\
E_\theta &= \delta_z \frac{i \beta}{k_c ^ 2} \frac{n}{r} J_n(r k_c)\sin(n \theta) \\
H_r &= -\frac{i \epsilon \omega}{k_c ^ 2} \frac{n}{r} J_n(r k_c)\sin(n \theta) \\
H_\theta &= -\frac{i \epsilon \omega}{k_c ^ 2} k_c J'_n(r k_c)\cos(n \theta)
\end{align}

Soit dans le plan transvere:

\begin{align}
\vec{E_t} = \delta_z \frac{i \beta}{k_c ^ 2} \begin{pmatrix} -k_c J'_n(r k_c)\cos(n \theta) \\ \frac{n}{r} J_n(r k_c)\sin(n \theta) \end{pmatrix} e^{-i\delta_z\beta z} \\
\vec{H_t} = \frac{i \epsilon \omega}{k_c ^ 2} \begin{pmatrix} -\frac{n}{r} J_n(r k_c)\sin(n \theta) \\ -k_c J'_n(r k_c)\cos(n \theta) \end{pmatrix} e^{-i\delta_z\beta z}
\end{align}

Dans la suite de ce rapport, on notera les champs sous la forme compacte suivante:

\begin{align}
\vec{E_t} = E_0 \vec{e}(r, \theta) e^{-i\delta_z\beta z} \\
\vec{H_t} = \delta_z Y E_0 \vec{h}(r, \theta) e^{-i\delta_z\beta z}
\end{align}

avec Y l'admittance du mode: $\frac{\beta}{\mu \omega}$ pour un mode TE, $\frac{\epsilon \omega}{\beta}$ pour un mode TM

% ==========================================================
\chapter{Calcul des champs à l'interface entre deux guides d'onde circulaires de rayons différents}
% ==========================================================

\section{Description géométrique du problème}
On considère la jonction à l'origine $z = 0$ entre deux guides d'onde circulaires alignés sur le même axe :
\begin{itemize}
    \item \textbf{Guide 1} ($z < 0$) : rayon $a$.
    \item \textbf{Guide 2} ($z > 0$) : rayon $b$.
\end{itemize}
On suppose sans perte de généralité que $a < b$ (discontinuité par élargissement). À l'interface $z = 0$, $S_1$ désigne la section du guide 1 et $S_2$ celle du guide 2. La région $S_2 \setminus S_1$ correspond au masque métallique conductrice parfaite (discontinuité de rayon).

\section{Décomposition modale des champs}
La méthode employée pour résoudre cette discontinuité est la \textbf{Technique de Raccordement de Modes} (\textit{Mode Matching Technique}).

Dans le \textbf{Guide 1} ($z \le 0$), le champ transversal total est la somme du champ incident et des modes réfléchis :
\begin{align}
\vec{E}_t^{(1)}(r,\theta,z) &= \sum_{j=1}^{M} \left( {a_j}^{+} e^{-i\beta_j^{(1)} z} + {a_j}^{-} e^{+i\beta_j^{(1)} z} \right) \vec{e}_j^{(1)}(r,\theta) \\
\vec{H}_t^{(1)}(r,\theta,z) &= \sum_{j=1}^{M} Y_j^{(1)} \left( {a_j}^{+} e^{-i\beta_j^{(1)} z} - {a_j}^{-} e^{+i\beta_j^{(1)} z} \right) \vec{h}_j^{(1)}(r,\theta)
\end{align}

De même dans le \textbf{Guide 2} ($z \ge 0$):
\begin{align}
\vec{E}_t^{(2)}(r,\theta,z) &= \sum_{k=1}^{N} \left( {b_k}^{+} e^{-i\beta_k^{(2)} z} + {b_k}^{-} e^{+i\beta_k^{(2)} z} \right) \vec{e}_k^{(2)}(r,\theta) \\
\vec{H}_t^{(2)}(r,\theta,z) &= \sum_{k=1}^{N} Y_k^{(2)} \left( {b_k}^{+} e^{-i\beta_k^{(2)} z} - {b_k}^{-} e^{+i\beta_k^{(2)} z} \right) \vec{h}_k^{(2)}(r,\theta)
\end{align}

où ${a_i}^{+}, {a_i}^{-}$ sont respectivement les amplitudes modales incidentes et réfléchies dans le guide 1, ${b_k}^{+}, {b_k}^{-}$ les amplitudes des modes réfléchis et incidents dans le guide 2, et $Y^{(1,2)}$ les admittances modales.

\section{Conditions de continuité à l'interface ($z = 0$)}
Au niveau de la discontinuité $z = 0$, les équations aux limites s'écrivent :

1. Continuité du champ électrique transversal $\vec{E}_t$ :
\begin{equation}
\begin{cases} 
\vec{E}_t^{(1)}(r,\theta,0) = \vec{E}_t^{(2)}(r,\theta,0) & \text{sur la surface d'ouverture } S_1 (r \le a) \\
\vec{E}_t^{(2)}(r,\theta,0) = 0 & \text{sur la paroi métallique } S_2 \setminus S_1 (a < r \le b)
\end{cases}
\end{equation}

2. Continuité du champ magnétique transversal $\vec{H}_t$ sur l'ouverture $S_1$ :
\begin{equation}
\vec{H}_t^{(1)}(r,\theta,0) = \vec{H}_t^{(2)}(r,\theta,0) \quad \text{pour } r \le a
\end{equation}

\section{Formulation matricielle et intégrales de recouvrement}
En appliquant les propriétés d'orthogonalité des modes propres :
\begin{align}
\iint_{S} \vec{e}_i \cdot \vec{e}_j \, dS = \delta_{ij} \\
\iint_{S} \vec{h}_i \cdot \vec{h}_j \, dS = \delta_{ij}
\end{align}

On projette la continuité de $\vec{E}_t$ sur les modes du guide 2 ($\vec{e}_n^{(2)}$) sur la surface $S_1$ :
\begin{equation}
\sum_{j=1}^{M} \left( {a_j}^{+} + {a_j}^{-} \right) \iint_{S_1} \vec{e}_j^{(1)} \cdot \vec{e}_n^{(2)} \, dS = {b_n}^{+} + {b_n}^{-}
\end{equation}

Soit sous forme vectorielle/matricielle :
\begin{equation}
\mathbf{M}^T (\mathbf{a^+} + \mathbf{a^-}) = \mathbf{b^+} + \mathbf{b^-} 
\label{eq:modal_E}
\end{equation}
où $\mathbf{M}$ est la matrice d'intercouplage modale de dimensions $M \times N$, dont les éléments sont définis par les intégrales de recouvrement sur la surface commune $S_1$ :
\begin{equation}
M_{mn} = \iint_{S_1} \vec{e}_m^{(1)}(r,\theta) \cdot \vec{e}_n^{(2)}(r,\theta) \, r \, dr \, d\theta
\end{equation}

De même, en projetant la continuité du champ magnétique $\vec{H}_t$ sur les modes du guide 1 ($\vec{h}_m^{(1)}$) sur la surface $S_1$ :
\begin{equation}
Y_m^{(1)} \left( {a_m}^{+} - {a_m}^{-}\right) = \sum_{k=1}^{N} Y_k^{(2)} \left( {b_k}^{+} - {b_k}^{-} \right) \iint_{S_1} \vec{h}_m^{(1)} \cdot \vec{h}_k^{(2)} \, dS
\end{equation}

Ce qui donne, sous forme matricielle:

\begin{equation}
\mathbf{Y}^{(1)} (\mathbf{a^+} - \mathbf{a^-}) = \mathbf{M} \mathbf{Y}^{(2)} (\mathbf{b^+} - \mathbf{b^-})
\label{eq:modal_H}
\end{equation}
où $\mathbf{Y}^{(1)}$ et $\mathbf{Y}^{(2)}$ sont les matrices diagonales d'admittances modales.

\section{Résolution de la matrice de répartition (Matrice S)}

On cherche a calculer la matrice S au niveau de la discontinuité:

\begin{equation}
\begin{bmatrix} \mathbf{a}^- \\ \mathbf{b}^+ \end{bmatrix} = \begin{bmatrix} \mathbf{S}_{11} & \mathbf{S}_{12} \\ \mathbf{S}_{21} & \mathbf{S}_{22} \end{bmatrix} \begin{bmatrix} \mathbf{a}^+ \\ \mathbf{b}^- \end{bmatrix}
\end{equation}

On isole le terme $\mathbf{b^+}$ dans l'équation \eqref{eq:modal_E}:

\begin{equation}
\label{eq:expr_b+}
\mathbf{b}^+ = \mathbf{M}^T \mathbf{a}^+ + \mathbf{M}^T \mathbf{a}^- - \mathbf{b}^-
\end{equation}

Et on le ré-injecte dans l'équation \eqref{eq:modal_H}:

\begin{equation}
\mathbf{Y}^{(1)} (\mathbf{a}^+ - \mathbf{a}^-) = \mathbf{M} \mathbf{Y}^{(2)} \left( [\mathbf{M}^T \mathbf{a}^+ + \mathbf{M}^T \mathbf{a}^- - \mathbf{b}^-] - \mathbf{b}^- \right)
\end{equation}

qui donne après simplification:

\begin{equation}
\mathbf{Y}^{(1)} \mathbf{a}^+ - \mathbf{Y}^{(1)} \mathbf{a}^- = \mathbf{M} \mathbf{Y}^{(2)} \mathbf{M}^T \mathbf{a}^+ + \mathbf{M} \mathbf{Y}^{(2)} \mathbf{M}^T \mathbf{a}^- - 2 \mathbf{M} \mathbf{Y}^{(2)} \mathbf{b}^-
\end{equation}

En regroupant les termes en $\mathbf{a}^-$ à gauche et les termes en $\mathbf{a}^+$ et $\mathbf{b}^-$ à droite :

\begin{equation}
(\mathbf{Y}^{(1)} + \mathbf{M} \mathbf{Y}^{(2)} \mathbf{M}^T) \mathbf{a}^- = (\mathbf{Y}^{(1)} - \mathbf{M} \mathbf{Y}^{(2)} \mathbf{M}^T) \mathbf{a}^+ + 2 \mathbf{M} \mathbf{Y}^{(2)} \mathbf{b}^-
\end{equation}

En posant $\mathbf{A} = \mathbf{Y}_1 + \mathbf{M} \mathbf{Y}_2 \mathbf{M}^T$:

\begin{equation}
\mathbf{a}^- = \mathbf{A}^{-1} (\mathbf{Y}^{(1)} - \mathbf{M} \mathbf{Y}^{(2)} \mathbf{M}^T) \mathbf{a}^+ + 2 \mathbf{A}^{-1} \mathbf{M} \mathbf{Y}^{(2)} \mathbf{b}^-
\end{equation}

ce qui permet de déterminer les expressions des blocs $\mathbf{S}_{11}$ et $\mathbf{S}_{12}$:

\begin{align}
\mathbf{S}_{11} = (\mathbf{Y}^{(1)} + \mathbf{M} \mathbf{Y}^{(2)} \mathbf{M}^T)^{-1} (\mathbf{Y}^{(1)} - \mathbf{M} \mathbf{Y}^{(2)} \mathbf{M}^T) \\
\mathbf{S}_{12} = 2 (\mathbf{Y}^{(1)} + \mathbf{M} \mathbf{Y}^{(2)} \mathbf{M}^T)^{-1} \mathbf{M} \mathbf{Y}^{(2)}
\end{align}

On réinjecte maintenant l'expression obtenue pour $\mathbf{a}^-$ dans l'équation \eqref{eq:expr_b+}:

\begin{equation}
\mathbf{b}^+ = \mathbf{M}^T \mathbf{a}^+ + \mathbf{M}^T (\mathbf{S}_{11} \mathbf{a}^+ + \mathbf{S}_{12} \mathbf{b}^-) - \mathbf{b}^-
\end{equation}

après simplification:
\begin{equation}
\mathbf{b}^+ = \mathbf{M}^T (\mathbf{I}_N + \mathbf{S}_{11}) \mathbf{a}^+ + (\mathbf{M}^T \mathbf{S}_{12} - \mathbf{I}_M) \mathbf{b}^-
\end{equation}

qui donne les expressions de $\mathbf{S}_{21}$ et $\mathbf{S}_{22}$:

\begin{align}
\mathbf{S}_{21} = 2 \mathbf{M}^T (\mathbf{Y}^{(1)} + \mathbf{M} \mathbf{Y}^{(2)} \mathbf{M}^T)^{-1} \mathbf{Y}^{(1)} \\
\mathbf{S}_{22} = 2 \mathbf{M}^T (\mathbf{Y}^{(1)} + \mathbf{M} \mathbf{Y}^{(2)} \mathbf{M}^T)^{-1} \mathbf{M} \mathbf{Y}^{(2)} - \mathbf{I}_M
\end{align}

\end{document}